\section{Team Roles and Communication}
\label{sec:team_roles}

\subsection{Role Allocation}

Roles were assigned based on individual expertise, interest, and prior experience. We aimed for each member to have a clearly defined primary responsibility while maintaining enough cross-training that no single absence would block progress~\cite{katzenbach2015wisdom}. The team naturally formed around a T-shaped competency model: each member had deep expertise in one area (the vertical bar of the T) and broad understanding of the overall system (the horizontal bar).

\subsection{RACI Matrix}

We used a RACI (Responsible, Accountable, Consulted, Informed) matrix to clarify ownership of key deliverables and avoid ambiguity~\cite{pmbok2021}. \tabref{tab:raci} presents the matrix for major project activities.

\begin{table}[htbp]
\centering
\compactTable
\caption{RACI matrix for key project activities. R = Responsible, A = Accountable, C = Consulted, I = Informed.}
\label{tab:raci}
\begin{tabular}{l c c c c c}
\toprule
\textbf{Activity} & \textbf{Dima} & \textbf{Robin} & \textbf{Pilot} & \textbf{Hardware} & \textbf{PM} \\
\midrule
CV model development & R/A & I & I & I & I \\
State machine design & C & R/A & I & I & I \\
Pi software integration & R/A & C & I & C & I \\
Hardware assembly \& wiring & C & C & I & R/A & I \\
Flight controller config & I & C & C & R/A & I \\
RC setup \& kill switch & I & I & R/A & C & I \\
Flight test piloting & I & C & R/A & I & I \\
Ground station operation & C & R/A & I & I & I \\
Risk assessment & C & C & C & C & R/A \\
Report writing & C & C & C & C & R/A \\
Schedule management & I & I & I & I & R/A \\
Simulation development & R/A & C & I & I & I \\
Progressive test planning & R/A & C & C & C & C \\
Search area survey & I & I & I & I & R/A \\
\bottomrule
\end{tabular}
\end{table}

\subsection{Communication Plan}

Effective communication was critical given the interdisciplinary nature of the project and the need to coordinate hardware availability, software integration, and field testing logistics. We established the following communication channels:

\begin{table}[htbp]
\centering
\compactTable
\caption{Communication channels and their purposes.}
\label{tab:comms}
\begin{tabularx}{\linewidth}{l l X}
\toprule
\textbf{Channel} & \textbf{Frequency} & \textbf{Purpose} \\
\midrule
Group chat & Daily & Quick questions, sharing results, coordination \\
Weekly meeting & Weekly & Sprint planning, progress review, design decisions \\
GitHub pull requests & As needed & Code review, integration approval \\
\texttt{GROUP\_STATUS.md} & Continuous & Persistent record of decisions, blockers, action items \\
Ad-hoc pair sessions & As needed & Debugging hardware issues, integration testing \\
\bottomrule
\end{tabularx}
\end{table}

\subsubsection{Decision Documentation}

All significant technical decisions were recorded in a \texttt{DESIGN\_DECISIONS.md} file using a structured format: context, options considered, decision, rationale, and trade-offs accepted. This served two purposes: (1) it prevented re-litigating settled decisions in later meetings, and (2) it provided ready-made content for the project report. Over the project duration, we documented ten formal design decisions (DD-01 through DD-10) covering topics from streaming protocols to geofence strategies.

\subsection{Team Dynamics}

The team progressed through Tuckman's stages of group development~\cite{tuckman1965developmental} in an accelerated fashion due to the short project timeline:

\begin{itemize}
  \item \textbf{Forming} (weeks 1--2): initial meetings, getting to know each other's skills, agreeing on tools and workflow.
  \item \textbf{Storming} (weeks 2--3): negotiating responsibilities, resolving differences in coding style and development workflow (e.g.\ branch naming, commit conventions).
  \item \textbf{Norming} (weeks 3--5): established routines, clear ownership, effective use of GitHub and group chat.
  \item \textbf{Performing} (weeks 5--12): focused execution, parallel workstreams running smoothly, integration checkpoints well-attended.
\end{itemize}

A key enabler of smooth collaboration was the \textbf{modular software architecture}. Because each major component (\texttt{vision.py}, \texttt{planning.py}, \texttt{config.py}, \texttt{main.py}) had well-defined interfaces, team members could work on their respective areas without creating merge conflicts or breaking each other's code. The vision system, for example, exposed a single function --- \texttt{detect\_in\_image(frame)} returning \texttt{(found, x, y, conf)} --- that all other modules consumed. This interface contract meant the CV developer could change model backends, preprocessing, or inference logic without touching any other file.

\subsection{Knowledge Transfer}

Given the specialised nature of certain components (particularly the AI pipeline and MAVLink protocol), we invested in knowledge transfer activities:

\begin{itemize}
  \item \textbf{Documentation}: comprehensive \texttt{CLAUDE.md} (project ground truth), architecture documents, setup guides for each platform, and flight day checklists --- all version-controlled alongside the code.
  \item \textbf{Pair programming}: critical integration points (e.g.\ connecting the Pi to the Cube flight controller via MAVProxy) were done in pairs to spread knowledge.
  \item \textbf{Written guides}: step-by-step setup documents (\texttt{PI\_SETUP.md}, \texttt{FLIGHT\_DAY\_CHECKLIST.md}) ensured any team member could reproduce the environment independently.
  \item \textbf{Interactive simulator}: \texttt{simple\_simulator.py} allowed any team member to experience the full mission flow on their laptop without hardware, building shared understanding of the state machine and detection logic.
\end{itemize}
